% !TeX root = 20261013.tex
% !LW recipe=lablog XeLaTeX
\input{_support/style.tex}

\begin{document}

\reporttitle{2026-10-13}{실험 결과 및 설계 기록}

\section*{목적과 이 보고서의 범위}

이 날짜의 실험 설정, 관찰 결과, 원인에 대한 해석과 후속 판단을 함께 정리한다. 다른 날짜의 평가 방식이나 표본 수가 다른 점수를 직접 합치지 않으며, 실행 완료와 논문 근거 채택을 구별한다.

\recordnote{수합·문서 개편: 2026-10-20. 실행 날짜는 각 절과 원시 실행 목록을 따른다. 아래의 과거 계획·진행 상태는 당시 기록이며, 현재 채택 상태와 최종 결과는 후속 실험을 함께 확인한다.}

\section{EXP71 · 같은 공급 사례의 누적 전체 비교}
2026-10-13 사용자가 제한된 과거 결과를 검토한 뒤 “keep going”으로 후속 실행을 요청했다. 승인된 연구 버전 v7 /\allowbreak{} EXP71 안에서 누적 전체 조건을 별도 run으로 실행한다. 기존 제한된 과거 결과는 보존한다. 10만 행 공급 확대는 이번 누적 비교 이후의 단계로 남긴다.

기준 실행은 \texttt{tabpfn/\allowbreak{}\allowbreak{}r\allowbreak{}esults/\allowbreak{}\allowbreak{}v\allowbreak{}7/\allowbreak{}\allowbreak{}exp71/\allowbreak{}\allowbreak{}20261013\allowbreak{}\_\allowbreak{}\allowbreak{}140543\_\allowbreak{}\allowbreak{}bounded\_\allowbreak{}\allowbreak{}uniform\_\allowbreak{}\allowbreak{}natural\_\allowbreak{}\allowbreak{}s43\_\allowbreak{}\allowbreak{}731a\allowbreak{}4a80}이다. 이 실행의 준비된 train bank, 단계별 신규 공급 IDs, 정제 test 전체를 그대로 사용한다. 공격 균등 /\allowbreak{} 공격 내 자연비율, Benign 75\%의 단계별 공급, 최초 출현 순서, seed 43, 두 모델의 설정을 유지한다. 미래 클래스의 사례는 최초 도입 전에 사용하지 않는다.

각 단계에서 그때까지 공급된 모든 행을 학습 집합으로 사용한다. XGBoost는 해당 집합으로 새로 학습하고 TabPFN은 같은 집합으로 전처리 및 context 캐시를 재구성한다. 이전 단계에서 버리는 행이 없으며, 최종 입력은 CIC2018 22,280행, ToN 36,864행이다. CLI의 \texttt{-\allowbreak{}-\allowbreak{}histor\allowbreak{}y-\allowbreak{}cap 40\allowbreak{}96}은 공급 stream의 기존 설정 일치 확인용이며 \texttt{-\allowbreak{}-\allowbreak{}histor\allowbreak{}y-\allowbreak{}mode c\allowbreak{}umulativ\allowbreak{}e}에서는 보존 상한으로 작동하지 않는다.

이번 비교는 과거 보존 정책을 바꾼 결과다. 총 입력량뿐 아니라 클래스별 잔존 수와 비율도 변하므로 효과를 데이터 양 하나로 설명하지 않는다. 특히 ToN 자연비율의 신규 MitM은 양쪽 모두 같은 11행, Ransomware 최초 공급은 같은 8행이다. 마지막 단계의 Ransomware 보존은 제한 조건 1행과 누적 조건 8행으로 달라진다.

갱신 시간은 동일한 GPU 동기화 구간에서 입력 구성·모델 생성·fit/\allowbreak{}cache의 합을 측정한다. 초기 구축은 별도이며 test 추론과 결과 기록·배포 비용을 갱신 시간에 합산하지 않는다. GPU 모델 worker는 순차 실행해 비용 경합을 피한다. Seed는 단일이며 비용은 단계별 1회 관측값이다. 실제 규모에 따른 비용 차이는 측정 후 판단한다.

두 모델 및 두 보존 정책 사이에 test IDs가 일치하는지, 누적 조건에서 모든 과거 공급 IDs가 남았는지, 각 단계의 신규 공급이 기준 실행과 같은지 검증한다. 저장한 확률·예측으로 클래스별 P/\allowbreak{}R/\allowbreak{}F1, 정상 오탐, 동일 기존 test 행의 유지·교정·훼손을 재계산한다. 완료 후 10/\allowbreak{}13 HTML 탭에 제한/\allowbreak{}누적 비교와 소표본 공격의 해석을 추가하고 연구 기록을 남긴다.

2026-10-13 시스템 재부팅으로 첫 누적 실행이 중단되었다. 사용자가 NVIDIA 복구 후 재개를 요청하여 \texttt{recover\_\allowbreak{}\allowbreak{}exp71.\allowbreak{}py\allowbreak{}}로 같은 v7/\allowbreak{}EXP71 안에 별도 복구 실행을 등록했다. 첫 실행과 손상 파일은 덮어쓰지 않는다. CIC2018 두 배정의 두 모델 및 ToN 균등의 두 모델은 완료된 32단계를 복사하고, ToN 자연비율은 두 모델 모두 초기 단계부터 다시 실행한다. TabPFN의 초기 cold fit와 후속 warm fit 측정이 섞이지 않도록 부분 worker에서 바로 이어 측정하지 않는다.

복구 후 커널은 7.0.0-38, NVIDIA 드라이버는 580.178.04이다. 기존 제한 조건과 복사한 누적 결과는 커널 7.0.0-31, 드라이버 580.173.02에서 측정했다. GPU 모델·Python 패키지·TabPFN 소스·checkpoint의 일치를 검사하고, worker별 환경 및 재사용 출처를 \texttt{recovery\allowbreak{}.\allowbreak{}json}에 기록한다. ToN 자연비율의 제한/\allowbreak{}누적 비용 차이는 환경 변경도 포함하므로 보존 정책만의 인과 효과로 해석하지 않는다. 재실행한 초기 두 단계의 예측이 기존 완료 결과와 일치하는지도 별도로 검증한다.



\section{EXP71 · 제한된 과거 사례로 새 공격 반영}
2026-10-13 사용자 실행·HTML 보고 승인. 기존 승인 연구 버전 v7에서 실행한다. 신규 연구 버전은 만들지 않는다. 이번 실행은 제한된 과거 조건만 다루고, 누적 전체는 결과 검토 후 같은 stream IDs로 별도 실행한다. Expert·S/\allowbreak{}V는 제외한다.

두 데이터셋 CIC2018·ToN, seed 43, 단일 TabPFN-v3와 XGBoost를 비교한다. 공격의 최초 출현 순서는 EXP69/\allowbreak{}70에서 검증한 순서를 유지한다. Benign은 모든 단계에 새 사례를 공급한다. Raw timestamp 재현이나 네트워크 환경 변화 실험은 아니다.

Case 1은 각 공격 1,024행을 배정하되 CIC2018 Web attack만 가용 train 450행으로 제한한다. Case 2는 Case 1의 전체 공격 예산(CIC2018 5,570, ToN 9,216)을 정제 train의 공격 클래스별 빈도 비례로 정수 배분한다. 자연 분포는 공격끼리의 비율이며 Benign을 포함한 전체 자연 분포를 재현하지 않는다. Test 빈도는 배분에 쓰지 않는다.

최종 Benign 75\%를 맞춰 CIC2018 16,710행, ToN 27,648행을 공급한다. CIC2018은 초기 포함 5단계에 각 3,342행, ToN은 6단계에 각 4,608행이다. Benign IDs와 단계별 공급은 두 case에서 동일하다. 최종 공급 총량은 각각 22,280행, 36,864행이다. 공격별 공급량과 시점이 case 간 달라지므로 단계별 업데이트 규모도 함께 보고한다.

매 단계 직전 학습 집합에서 최대 4,096행을 보존하고 이번 단계의 신규 Benign·공격 사례를 모두 더한다. 따라서 4,096은 전체 context 상한이 아니라 과거 보존 상한이다. 이미 버린 과거 사례는 다시 가져오지 않는다. 과거 선택은 Benign 목표 75\%와 해당 case의 공격 비율을 사용하며, 직전 집합의 클래스별 가용량을 넘을 수 없다. 이미 도입한 클래스는 최소 1행을 유지한다. 최소 보존 때문에 자연 비율이 보정된 클래스와 보정 전후 quota를 기록한다. 가용량 부족 시 전체 보존 예산을 다른 클래스로 재배분한다.

두 모델에는 같은 학습·평가 IDs를 제공하고 미래 클래스가 context·전처리에 들어가지 않는지 검증한다. 모델 설정은 EXP69/\allowbreak{}70과 같다: TabPFN v3\_\allowbreak{}20260417\_\allowbreak{}multiclass, ensemble 4, frozen backbone, fit\_\allowbreak{}with\_\allowbreak{}cache, CPU cache; XGB 300 rounds, depth 8, lr .05, CUDA hist. GPU worker는 하나씩 실행해 비용 경합을 피한다. 초기 구축과 후속 갱신을 구분한다.

갱신 시간은 과거 선택·입력 구성, 모델 생성, XGB 학습 또는 TabPFN 전처리·context 캐시 구성의 GPU 동기화된 경과 시간 합이다. 최초 checkpoint 적재는 초기 fit에 포함될 수 있다. 갱신된 모델이 메모리에서 준비되는 시점까지 측정하며 결과 파일 기록, 운영 배포·모델 직렬화, 전체 test 추론은 포함하지 않는다. 첫 추론 batch 지연과 전체 test 추론 시간은 별도로 저장한다. Test 전처리는 기존 정제 기준을 따르고 별도 tuning은 수행하지 않는다.

평가는 기존 정제 test 중 도입한 클래스의 전체 행을 사용한다. 클래스별 P/\allowbreak{}R/\allowbreak{}F1·support와 Benign 오탐을 저장한다. 이전 클래스 유지 평가는 새 클래스의 정답 행만 제외하고 새 클래스로의 오분류는 유지한 동일 과거 표본군에서 계산한다. 갱신 전후 교정·훼손도 같은 행에서 비교한다. 동일 feature가 context에 있는 test 행을 제외한 보조 평가를 제공한다. 기존 개발 holdout·단일 seed라는 범위는 유지한다.

재현을 위해 단계별 유입 IDs, 보존 IDs, 학습 IDs, 평가 IDs, 확률, 예측, confusion matrix, 비용 phase, checkpoint·입력·코드 hash를 보존한다. 준비된 공급 stream을 나중에 \texttt{-\allowbreak{}-\allowbreak{}histor\allowbreak{}y-\allowbreak{}mode c\allowbreak{}umulativ\allowbreak{}e -\allowbreak{}-\allowbreak{}stre\allowbreak{}am-\allowbreak{}root \allowbreak{}EXISTING\allowbreak{}\_\allowbreak{}\allowbreak{}RUN}으로 재사용할 수 있다. 실행 성공과 논문 주장 채택은 구별한다.



\section{2026-10-13 · EXP71 · 새 공격 도입 시나리오에서 TabPFN을 채택할 근거}
\label{report:20261013_exp71_cumulative_results}
\recordnote{이 절은 당시 설정의 결과다. 평가 방식 변경·정정·후속 비교는 뒤에 명시한 실험과 구별한다.}
2026-10-14 보고서 재구성. 10/\allowbreak{}13에 완료한 제한 보존·누적 전체 실험을 사용했으며 추가 추론이나 새 실험은 수행하지 않았다.

연구 의도는 새 공격의 라벨 사례가 확보됐을 때 사전학습 가중치를 고정한 TabPFN이 context 변경만으로 유용한 탐지 품질을 유지하면서 갱신 비용을 줄일 수 있는지 확인하는 것이다. 원본 클래스의 최초 출현 순서를 유지하되 초기 및 일부 인접 공격을 묶었고, 각 단계에서 신규 공격과 Benign 사례를 공급했다. 실제 시간축의 트래픽 재현이나 미관측 공격의 라벨 없는 탐지를 검증한 것은 아니다.

변화 조건은 공격별 공급량(균등 /\allowbreak{} 공격 내 자연 비율)과 과거 보존(최대 4,096개 /\allowbreak{} 누적 전체)이다. 두 데이터에서 8개 구성을 비교했다. 최종 공급은 CIC2018 22,280행, ToN 36,864행이고 Benign은 각 75\%를 매 단계 나눠 공급한다. 구성 내 두 모델의 학습·test IDs와 보존 정책 간 신규 공급은 고정했다.

제한 보존에서는 4개 구성 중 3개에서 TabPFN의 후속 갱신 합계가 더 짧았고, 누적 전체에서는 0개였다. 최종 test의 대량 추론은 8개 구성 모두 XGBoost가 더 빨랐다. 따라서 갱신 비용 절감과 추론 비용 절감을 구별해야 한다.

갱신 시간은 초기 구축을 제외한 후속 단계 합계다. 추론은 비교 범위를 맞추기 위해 최종 단계의 전체 test 1회 시간을 별도로 제시한다. 배치 예측과 결과 기록이 포함되고, 운영 배포·직렬화는 미측정이다. 두 열의 측정 범위가 다르므로 단순 합계를 운영 총비용으로 해석하지 않는다.

CIC2018 균등·제한 보존에서는 두 모델의 최종 Macro-F1이 약 76.12\%로 유사했다. 같은 균등 공급의 누적 조건에서는 TabPFN 79.13\% /\allowbreak{} XGB 77.60\%였다. CIC2018 자연 배정과 ToN의 두 배정에서는 최종 Macro-F1이 XGBoost보다 낮아, 데이터·구성 전반의 우위로 이어지지는 않았다.

채택을 검토할 근거 ① — CIC2018 균등·제한 보존: 후속 갱신 합계는 TabPFN 2.74초 /\allowbreak{} XGB 3.67초, 최종 F1은 76.11\% /\allowbreak{} 76.12\%였다. 정상 오탐률은 0.0600\% /\allowbreak{} 1.1735\%다. 비슷한 전체 F1에서 갱신과 오탐의 이점이 함께 관찰됐다. 다만 최종 test 추론은 TabPFN이 더 느려 운영 총비용 우위를 뜻하지 않는다.

채택을 검토할 근거 ② — CIC2018 균등·누적 전체: Infiltration F1은 TabPFN 41.31\% /\allowbreak{} XGB 11.37\%로 특정 소수 공격에서 이점이 있었다. 그러나 이 조건의 TabPFN 갱신은 더 느렸고 Web attack까지 같은 이점이 나타난 것은 아니다. 공격별 탐지 품질을 선택적으로 검증할 이유다.

아직 해결할 점 — ToN 자연비율·누적 전체: TabPFN 최종 F1 69.30\%는 XGB 76.80\%보다 낮다. TabPFN의 제한 \ensuremath{\rightarrow} 누적 개선에는 Ransomware F1 0.00\% \ensuremath{\rightarrow} 36.54\%가 포함되지만, Scanning F1은 0.00\%이고 MitM precision은 4.08\%에 머문다. 과거 사례를 더 남기는 것만으로 tail 문제가 해결되지는 않았다.

현재의 판단 — 가중치 재학습 없이 새 클래스를 반영할 수 있다는 작동 방식과 제한된 조건의 이점은 확인했다. 모든 공격의 tail 성능 개선이나 XGBoost보다 낮은 운영 총비용은 아직 확인하지 못했다. TabPFN을 채택하려면 확보 라벨·context 예산·갱신 빈도·갱신당 추론량을 명시하고, 필요한 탐지 품질을 충족하는 조건에서 비교해야 한다.

후속 검증은 한 번 갱신한 뒤 처리할 트래픽 규모와 필요한 새 공격 탐지 품질을 먼저 정하고, 같은 라벨 공급·출현 순서에서 입력 예산을 변화시키는 것이다. 갱신·추론 시간을 반복 측정하고 새 공격의 P/\allowbreak{}R/\allowbreak{}F1, 기존 공격 유지, 정상 오탐을 함께 평가한다. 제한 보존 및 전체 누적 XGBoost와 비교해 TabPFN의 이점이 유지되는 범위를 확인한다. 라벨 공급 확대와 과거 보존 확대는 별도 요인으로 구별한다.

재개 환경: ToN 자연비율의 누적 조건은 재부팅 후 커널 7.0.0-38 /\allowbreak{} NVIDIA 580.178.04에서 두 모델을 처음부터 다시 측정했다. 제한 조건과 나머지 누적 조건은 커널 7.0.0-31 /\allowbreak{} NVIDIA 580.173.02에서 측정했다. GPU·Python 패키지·모델 소스·checkpoint는 같지만, ToN 자연비율의 제한 \ensuremath{\rightarrow} 누적 비용 차이를 보존 정책만의 효과로 해석하지 않는다.

단일 seed·단계별 1회 timing·반복 관찰한 개발 holdout이라는 범위에서 해석한다. 10만 행 공급 확대는 미실행이며, expert·S/\allowbreak{}V가 들어간 제안 모델의 기여를 이 결과로 주장하지 않는다.

\subsection{상세 수치}
CIC2018 · 균등: XGB의 최종 F1은 76.12\% \ensuremath{\rightarrow} 77.60\%, 후속 갱신 합계는 3.67초 \ensuremath{\rightarrow} 4.27초, 정상 오탐은 1.1735\% \ensuremath{\rightarrow} 1.8613\%다; TabPFN의 최종 F1은 76.11\% \ensuremath{\rightarrow} 79.13\%, 후속 갱신 합계는 2.74초 \ensuremath{\rightarrow} 4.96초, 정상 오탐은 0.0600\% \ensuremath{\rightarrow} 0.0788\%다.

CIC2018 · 공격 내 자연: XGB의 최종 F1은 72.52\% \ensuremath{\rightarrow} 75.23\%, 후속 갱신 합계는 2.77초 \ensuremath{\rightarrow} 3.19초, 정상 오탐은 0.1861\% \ensuremath{\rightarrow} 0.0966\%다; TabPFN의 최종 F1은 71.84\% \ensuremath{\rightarrow} 73.30\%, 후속 갱신 합계는 2.80초 \ensuremath{\rightarrow} 5.22초, 정상 오탐은 0.0032\% \ensuremath{\rightarrow} 0.0077\%다.

ToN · 균등: XGB의 최종 F1은 67.12\% \ensuremath{\rightarrow} 70.93\%, 후속 갱신 합계는 8.05초 \ensuremath{\rightarrow} 10.57초, 정상 오탐은 10.9470\% \ensuremath{\rightarrow} 11.2375\%다; TabPFN의 최종 F1은 65.98\% \ensuremath{\rightarrow} 68.61\%, 후속 갱신 합계는 4.08초 \ensuremath{\rightarrow} 11.10초, 정상 오탐은 13.9460\% \ensuremath{\rightarrow} 13.5476\%다.

ToN · 공격 내 자연: XGB의 최종 F1은 70.86\% \ensuremath{\rightarrow} 76.80\%, 후속 갱신 합계는 6.67초 \ensuremath{\rightarrow} 8.86초, 정상 오탐은 9.5934\% \ensuremath{\rightarrow} 10.0749\%다; TabPFN의 최종 F1은 65.27\% \ensuremath{\rightarrow} 69.30\%, 후속 갱신 합계는 4.22초 \ensuremath{\rightarrow} 11.37초, 정상 오탐은 11.5618\% \ensuremath{\rightarrow} 12.9361\%다.

ToN 자연비율의 MitM: 최종 사용 사례는 제한 11개, 누적 11개다. XGB: F1 66.33\% \ensuremath{\rightarrow} 65.53\%, P 71.29\% \ensuremath{\rightarrow} 96.80\%, R 62.02\% \ensuremath{\rightarrow} 49.53\%, Benign \ensuremath{\rightarrow} MitM 205건 \ensuremath{\rightarrow} 16건; TabPFN: F1 5.74\% \ensuremath{\rightarrow} 7.64\%, P 3.01\% \ensuremath{\rightarrow} 4.08\%, R 61.15\% \ensuremath{\rightarrow} 59.78\%, Benign \ensuremath{\rightarrow} MitM 22,870건 \ensuremath{\rightarrow} 16,311건.

ToN 자연비율의 Ransomware: 최종 사용 사례는 제한 1개, 누적 8개다. XGB: F1 0.00\% \ensuremath{\rightarrow} 41.78\%, P 0.00\% \ensuremath{\rightarrow} 29.61\%, R 0.00\% \ensuremath{\rightarrow} 70.94\%, Benign \ensuremath{\rightarrow} Ransomware 0건 \ensuremath{\rightarrow} 1,196건; TabPFN: F1 0.00\% \ensuremath{\rightarrow} 36.54\%, P 0.00\% \ensuremath{\rightarrow} 35.61\%, R 0.00\% \ensuremath{\rightarrow} 37.52\%, Benign \ensuremath{\rightarrow} Ransomware 38건 \ensuremath{\rightarrow} 481건.

MitM의 신규 11개는 두 보존 정책에서 같다. 따라서 이 비교에서 관찰한 MitM 변화는 신규 MitM 라벨을 더 확보한 효과가 아니다. 다만 과거 정상·공격 사례의 수, 다양성 및 클래스 비율이 함께 바뀌므로 그중 하나를 단독 원인으로 확정하지 않는다. Ransomware는 마지막 잔존 수가 1개에서 8개로 달라져, 보존 정책에 따른 소수 클래스 유지 여부를 함께 확인할 수 있다.

\subsection{출처}
\par\noindent \textbullet\enspace 제한 실행: \texttt{tabpfn/\allowbreak{}\allowbreak{}r\allowbreak{}esults/\allowbreak{}\allowbreak{}v\allowbreak{}7/\allowbreak{}\allowbreak{}exp71/\allowbreak{}\allowbreak{}20261013\allowbreak{}\_\allowbreak{}\allowbreak{}140543\_\allowbreak{}\allowbreak{}bounded\_\allowbreak{}\allowbreak{}uniform\_\allowbreak{}\allowbreak{}natural\_\allowbreak{}\allowbreak{}s43\_\allowbreak{}\allowbreak{}731a\allowbreak{}4a80}


\par\noindent \textbullet\enspace 누적 실행: \texttt{tabpfn/\allowbreak{}\allowbreak{}r\allowbreak{}esults/\allowbreak{}\allowbreak{}v\allowbreak{}7/\allowbreak{}\allowbreak{}exp71/\allowbreak{}\allowbreak{}20261013\allowbreak{}\_\allowbreak{}\allowbreak{}174035\_\allowbreak{}\allowbreak{}cumulati\allowbreak{}ve\_\allowbreak{}\allowbreak{}recov\allowbreak{}ery\_\allowbreak{}\allowbreak{}s43\_\allowbreak{}\allowbreak{}4229fd93\allowbreak{}}


\par\noindent \textbullet\enspace 각 실행의 \texttt{summary.\allowbreak{}csv}, \texttt{class\_\allowbreak{}\allowbreak{}me\allowbreak{}trics.\allowbreak{}cs\allowbreak{}v}, \texttt{jobs/\allowbreak{}\allowbreak{}*/\allowbreak{}\allowbreak{}s\allowbreak{}tage\_\allowbreak{}\allowbreak{}*/\allowbreak{}\allowbreak{}m\allowbreak{}etrics.\allowbreak{}j\allowbreak{}son}


\par\noindent \textbullet\enspace 비교 검증: \texttt{tabpfn/\allowbreak{}\allowbreak{}r\allowbreak{}esults/\allowbreak{}\allowbreak{}v\allowbreak{}7/\allowbreak{}\allowbreak{}exp71/\allowbreak{}\allowbreak{}20261013\allowbreak{}\_\allowbreak{}\allowbreak{}174035\_\allowbreak{}\allowbreak{}cumulati\allowbreak{}ve\_\allowbreak{}\allowbreak{}recov\allowbreak{}ery\_\allowbreak{}\allowbreak{}s43\_\allowbreak{}\allowbreak{}4229fd93\allowbreak{}/\allowbreak{}\allowbreak{}compari\allowbreak{}son\_\allowbreak{}\allowbreak{}vali\allowbreak{}dation.\allowbreak{}j\allowbreak{}son}


\par\noindent \textbullet\enspace HTML: \texttt{docs/\allowbreak{}\allowbreak{}lab\allowbreak{}log/\allowbreak{}\allowbreak{}repo\allowbreak{}rts/\allowbreak{}\allowbreak{}2026\allowbreak{}1014.\allowbreak{}tex\allowbreak{}}




\section{2026-10-13 · EXP71 · 제한된 과거 사례의 새 공격 반영 결과}
\label{report:20261013_exp71_results}
\recordnote{이 절은 당시 설정의 결과다. 평가 방식 변경·정정·후속 비교는 뒤에 명시한 실험과 구별한다.}
2026-10-13. Seed 43, CIC2018·ToN, 균등/\allowbreak{}공격 내 자연 분포, 단일 TabPFN·XGB. 두 모델은 각 단계에서 동일한 학습·평가 IDs를 사용했다. 초기 구축을 제외한 갱신 시간을 중심으로 비교하며, 추론은 별도 기록한다. Expert/\allowbreak{}S/\allowbreak{}V 및 누적 전체 조건은 실행하지 않았다.

공급 라벨 전체를 모델에 계속 누적하지 않고 직전 학습 집합에서 최대 4,096행을 보존한 뒤 신규 사례를 더했다. 균등과 자연 배정은 최종 공격 총량과 Benign 공급을 공유한다. 자연 배정은 정제 train의 공격끼리의 비율이며, 적은 클래스의 보존에 최소 1행 조건을 적용했다. 이 조건의 클래스별 보정은 원본 결과에 남겼다.

CIC2018 /\allowbreak{} 균등에서 후속 갱신 시간 합계는 XGB 3.665초, TabPFN 2.735초로, TabPFN/\allowbreak{}XGB 비율은 0.746다. 최종 Macro-F1은 76.12\%와 76.11\%, Benign 오탐률은 1.1735\%와 0.0600\%다. 전체 22,280행을 공급했고 마지막 갱신에는 8,462행을 사용했다. 비용과 탐지 품질을 함께 판단해야 하며 단일 수치로 우위를 확정하지 않는다.

CIC2018 /\allowbreak{} 공격 내 자연에서 후속 갱신 시간 합계는 XGB 2.767초, TabPFN 2.796초로, TabPFN/\allowbreak{}XGB 비율은 1.010다. 최종 Macro-F1은 72.52\%와 71.84\%, Benign 오탐률은 0.1861\%와 0.0032\%다. 전체 22,280행을 공급했고 마지막 갱신에는 8,028행을 사용했다. 비용과 탐지 품질을 함께 판단해야 하며 단일 수치로 우위를 확정하지 않는다.

ToN /\allowbreak{} 균등에서 후속 갱신 시간 합계는 XGB 8.049초, TabPFN 4.083초로, TabPFN/\allowbreak{}XGB 비율은 0.507다. 최종 Macro-F1은 67.12\%와 65.98\%, Benign 오탐률은 10.9470\%와 13.9460\%다. 전체 36,864행을 공급했고 마지막 갱신에는 9,728행을 사용했다. 비용과 탐지 품질을 함께 판단해야 하며 단일 수치로 우위를 확정하지 않는다.

ToN /\allowbreak{} 공격 내 자연에서 후속 갱신 시간 합계는 XGB 6.670초, TabPFN 4.216초로, TabPFN/\allowbreak{}XGB 비율은 0.632다. 최종 Macro-F1은 70.86\%와 65.27\%, Benign 오탐률은 9.5934\%와 11.5618\%다. 전체 36,864행을 공급했고 마지막 갱신에는 8,715행을 사용했다. 비용과 탐지 품질을 함께 판단해야 하며 단일 수치로 우위를 확정하지 않는다.

\subsection{소표본 공격의 관찰}
CIC2018 · 균등의 평가 표본이 가장 적은 두 공격을 따로 보면, Infiltration는 test 8,742개, 최초 공급 1,024개, 최종 context 잔존 205개다. XGB F1 12.89\% (P 8.18\% /\allowbreak{} R 30.39\%, 공격 \ensuremath{\rightarrow} Benign 5,873건 /\allowbreak{} Benign \ensuremath{\rightarrow} 공격 29,812건); TabPFN F1 20.29\% (P 76.17\% /\allowbreak{} R 11.70\%, 공격 \ensuremath{\rightarrow} Benign 7,567건 /\allowbreak{} Benign \ensuremath{\rightarrow} 공격 320건). Web attack는 test 127개, 최초 공급 450개, 최종 context 잔존 204개다. XGB F1 24.08\% (P 13.83\% /\allowbreak{} R 92.91\%, 공격 \ensuremath{\rightarrow} Benign 8건 /\allowbreak{} Benign \ensuremath{\rightarrow} 공격 450건); TabPFN F1 13.88\% (P 7.49\% /\allowbreak{} R 93.70\%, 공격 \ensuremath{\rightarrow} Benign 8건 /\allowbreak{} Benign \ensuremath{\rightarrow} 공격 1,057건). 이 혼동 건수는 최종 예측에서 직접 확인한 관찰이며, 사례 수 부족이나 context 구성 중 무엇이 원인인지는 이 비교만으로 분리할 수 없다.

CIC2018 · 공격 내 자연의 평가 표본이 가장 적은 두 공격을 따로 보면, Infiltration는 test 8,742개, 최초 공급 143개, 최종 context 잔존 29개다. XGB F1 9.74\% (P 13.33\% /\allowbreak{} R 7.68\%, 공격 \ensuremath{\rightarrow} Benign 8,006건 /\allowbreak{} Benign \ensuremath{\rightarrow} 공격 4,364건); TabPFN F1 2.79\% (P 91.18\% /\allowbreak{} R 1.42\%, 공격 \ensuremath{\rightarrow} Benign 8,548건 /\allowbreak{} Benign \ensuremath{\rightarrow} 공격 12건). Web attack는 test 127개, 최초 공급 2개, 최종 context 잔존 1개다. XGB F1 0.00\% (P 0.00\% /\allowbreak{} R 0.00\%, 공격 \ensuremath{\rightarrow} Benign 127건 /\allowbreak{} Benign \ensuremath{\rightarrow} 공격 0건); TabPFN F1 1.49\% (P 14.29\% /\allowbreak{} R 0.79\%, 공격 \ensuremath{\rightarrow} Benign 125건 /\allowbreak{} Benign \ensuremath{\rightarrow} 공격 6건). 이 혼동 건수는 최종 예측에서 직접 확인한 관찰이며, 사례 수 부족이나 context 구성 중 무엇이 원인인지는 이 비교만으로 분리할 수 없다.

ToN · 균등의 평가 표본이 가장 적은 두 공격을 따로 보면, MitM는 test 1,161개, 최초 공급 1,024개, 최종 context 잔존 1,024개다. XGB F1 7.51\% (P 3.90\% /\allowbreak{} R 97.67\%, 공격 \ensuremath{\rightarrow} Benign 17건 /\allowbreak{} Benign \ensuremath{\rightarrow} 공격 20,722건); TabPFN F1 3.57\% (P 1.82\% /\allowbreak{} R 96.90\%, 공격 \ensuremath{\rightarrow} Benign 26건 /\allowbreak{} Benign \ensuremath{\rightarrow} 공격 50,428건). Ransomware는 test 709개, 최초 공급 1,024개, 최종 context 잔존 128개다. XGB F1 13.96\% (P 7.54\% /\allowbreak{} R 94.36\%, 공격 \ensuremath{\rightarrow} Benign 27건 /\allowbreak{} Benign \ensuremath{\rightarrow} 공격 8,204건); TabPFN F1 10.25\% (P 5.44\% /\allowbreak{} R 89.00\%, 공격 \ensuremath{\rightarrow} Benign 63건 /\allowbreak{} Benign \ensuremath{\rightarrow} 공격 10,919건). 이 혼동 건수는 최종 예측에서 직접 확인한 관찰이며, 사례 수 부족이나 context 구성 중 무엇이 원인인지는 이 비교만으로 분리할 수 없다.

ToN · 공격 내 자연의 평가 표본이 가장 적은 두 공격을 따로 보면, MitM는 test 1,161개, 최초 공급 11개, 최종 context 잔존 11개다. XGB F1 66.33\% (P 71.29\% /\allowbreak{} R 62.02\%, 공격 \ensuremath{\rightarrow} Benign 281건 /\allowbreak{} Benign \ensuremath{\rightarrow} 공격 205건); TabPFN F1 5.74\% (P 3.01\% /\allowbreak{} R 61.15\%, 공격 \ensuremath{\rightarrow} Benign 247건 /\allowbreak{} Benign \ensuremath{\rightarrow} 공격 22,870건). Ransomware는 test 709개, 최초 공급 8개, 최종 context 잔존 1개다. XGB F1 0.00\% (P 0.00\% /\allowbreak{} R 0.00\%, 공격 \ensuremath{\rightarrow} Benign 694건 /\allowbreak{} Benign \ensuremath{\rightarrow} 공격 0건); TabPFN F1 0.00\% (P 0.00\% /\allowbreak{} R 0.00\%, 공격 \ensuremath{\rightarrow} Benign 236건 /\allowbreak{} Benign \ensuremath{\rightarrow} 공격 38건). 이 혼동 건수는 최종 예측에서 직접 확인한 관찰이며, 사례 수 부족이나 context 구성 중 무엇이 원인인지는 이 비교만으로 분리할 수 없다.

관찰의 해석은 같은 단계의 같은 사례에서 두 갱신 방식이 보인 비용·예측 차이로 한정한다. 균등/\allowbreak{}자연 case 간에는 공격별 확보 라벨, 유입량과 실제 context 크기도 달라진다. 이를 분포만의 인과 효과로 설명하지 않는다. 특정 공격의 재현율 상승은 정상 및 다른 공격에서 유입된 FP와 함께 확인한다. 기존 클래스 유지 수치는 신규 정답 행을 제외한 고정 과거 표본군에서 계산하며 새 공격으로의 오분류는 오류에 포함한다.

갱신 시간은 메모리에서 입력 선택·모델 구성·fit/\allowbreak{}cache가 완료되기까지의 경과 시간이며 운영 배포·모델 직렬화를 측정하지 않았다. 첫 batch와 전체 test 추론 비용은 별도 지표다. 단일 seed·기존 개발 holdout·단계별 1회 timing이라는 범위에서 해석한다.

후속 판단은 같은 공급 IDs를 누적 전체 조건으로 재생하여 과거 보존 상한이 탐지 유지와 비용에 미친 영향을 비교하는 것이다. 이번 결과만으로 bounded 방식의 효과나 TabPFN 전체의 비용 우위를 일반화하지 않는다. 누적 전체 실행은 사용자의 결과 검토 후 진행한다.

\subsection{출처}
\par\noindent \textbullet\enspace 실행: \texttt{tabpfn/\allowbreak{}\allowbreak{}r\allowbreak{}esults/\allowbreak{}\allowbreak{}v\allowbreak{}7/\allowbreak{}\allowbreak{}exp71/\allowbreak{}\allowbreak{}20261013\allowbreak{}\_\allowbreak{}\allowbreak{}140543\_\allowbreak{}\allowbreak{}bounded\_\allowbreak{}\allowbreak{}uniform\_\allowbreak{}\allowbreak{}natural\_\allowbreak{}\allowbreak{}s43\_\allowbreak{}\allowbreak{}731a\allowbreak{}4a80}


\par\noindent \textbullet\enspace 단계별 수치: \texttt{tabpfn/\allowbreak{}\allowbreak{}r\allowbreak{}esults/\allowbreak{}\allowbreak{}v\allowbreak{}7/\allowbreak{}\allowbreak{}exp71/\allowbreak{}\allowbreak{}20261013\allowbreak{}\_\allowbreak{}\allowbreak{}140543\_\allowbreak{}\allowbreak{}bounded\_\allowbreak{}\allowbreak{}uniform\_\allowbreak{}\allowbreak{}natural\_\allowbreak{}\allowbreak{}s43\_\allowbreak{}\allowbreak{}731a\allowbreak{}4a80/\allowbreak{}\allowbreak{}sum\allowbreak{}mary.\allowbreak{}csv\allowbreak{}}


\par\noindent \textbullet\enspace 클래스별 수치: \texttt{tabpfn/\allowbreak{}\allowbreak{}r\allowbreak{}esults/\allowbreak{}\allowbreak{}v\allowbreak{}7/\allowbreak{}\allowbreak{}exp71/\allowbreak{}\allowbreak{}20261013\allowbreak{}\_\allowbreak{}\allowbreak{}140543\_\allowbreak{}\allowbreak{}bounded\_\allowbreak{}\allowbreak{}uniform\_\allowbreak{}\allowbreak{}natural\_\allowbreak{}\allowbreak{}s43\_\allowbreak{}\allowbreak{}731a\allowbreak{}4a80/\allowbreak{}\allowbreak{}cla\allowbreak{}ss\_\allowbreak{}\allowbreak{}metri\allowbreak{}cs.\allowbreak{}csv}


\par\noindent \textbullet\enspace 확률·예측·행 ID·혼동행렬: \texttt{tabpfn/\allowbreak{}\allowbreak{}r\allowbreak{}esults/\allowbreak{}\allowbreak{}v\allowbreak{}7/\allowbreak{}\allowbreak{}exp71/\allowbreak{}\allowbreak{}20261013\allowbreak{}\_\allowbreak{}\allowbreak{}140543\_\allowbreak{}\allowbreak{}bounded\_\allowbreak{}\allowbreak{}uniform\_\allowbreak{}\allowbreak{}natural\_\allowbreak{}\allowbreak{}s43\_\allowbreak{}\allowbreak{}731a\allowbreak{}4a80/\allowbreak{}\allowbreak{}job\allowbreak{}s/\allowbreak{}\allowbreak{}}


\par\noindent \textbullet\enspace 설계: \texttt{docs/\allowbreak{}\allowbreak{}lab\allowbreak{}log/\allowbreak{}\allowbreak{}repo\allowbreak{}rts/\allowbreak{}\allowbreak{}2026\allowbreak{}1013.\allowbreak{}tex\allowbreak{}}


\par\noindent \textbullet\enspace 보고서: \texttt{docs/\allowbreak{}\allowbreak{}lab\allowbreak{}log/\allowbreak{}\allowbreak{}repo\allowbreak{}rts/\allowbreak{}\allowbreak{}2026\allowbreak{}1014.\allowbreak{}tex\allowbreak{}}




\section*{출처와 실행 추적}

상세 수치의 원시 실행은 \path{inventory/runs.csv}, 실험별 연결은 \path{inventory/experiments.csv}에서 확인한다. 표·그림의 재현 입력은 \path{evidence/}와 \path{data/}에 둔다. 문서 이동·통합 내역은 \path{configs/experiments/migrations/20261020_documents/manifest.json}에 기록했다.

\end{document}
